\documentclass[11pt]{article}
\usepackage[margin=1in]{geometry}
\usepackage{amsmath,amssymb}
\usepackage[T1]{fontenc}
\usepackage[utf8]{inputenc}
\usepackage{lmodern}
\usepackage{hyperref}

\setlength{\parskip}{0.6em}
\setlength{\parindent}{0pt}

\title{Response to Reviewers}
\author{}
\date{}

\begin{document}
\maketitle

We thank the reviewers for the detailed second-round report. The revision responds to each remaining concern either by new technical content in the manuscript, by deliberate removal of underdeveloped material, or by an explicit acknowledgment where the honest answer is still ``open.'' This letter summarizes the global framing and then walks through the numbered concerns.

\section*{Global Framing}

We agree with the underlying standard the reviewers applied: if the paper merely restated the Gibbs variational identity, the MAP/MDL correspondence, and a deterministic Lindley-style information-gain criterion under new notation, the contribution would be largely synthetic. The second revision therefore had to do more than reposition the manuscript rhetorically.

The revised manuscript now distinguishes clearly between:
\begin{enumerate}
\item exact finite-history identities (the variational identity, the MAP/MDL bridge, and the risk-minus-information-gain decomposition),
\item a theorem-level interactive asymptotic layer (predictive merging, fixed-false-program elimination under cumulative R\'enyi-$\tfrac12$ separation, semantic-class posterior concentration under aggregate false-mixture separation, an explicit exploration-style sufficient condition, and an $\epsilon$-mixed expected-free-energy corollary), and
\item a conceptual evaluation framework built on top of that ideal theory: an ideal / behavioral proxy / mechanistic challenge three-level architecture with a concrete two-quantity compliance protocol.
\end{enumerate}

The paper's novelty claim is not that each ingredient is unprecedented in isolation. It is that the specific synthesis changes the latent object being inferred, the asymptotic target of learning, and the conditions under which active querying recovers explanation rather than merely prediction.

We have also removed the hard-to-vary material from the manuscript in response to RMC5 (see below). It was philosophically attractive but not adequately supported by theorems or by the compliance protocol, and keeping it weakened the paper's stance on what it actually delivers.

\section*{Response to RMC1: The Central Asymptotic Result Is Still a Conjecture}

The revision eliminates the conjecture. Section~8 now proves:
\begin{enumerate}
\item Predictive merging (Theorem): the universal interactive posterior predictive merges with the true completed kernel under any history-adapted policy, without any exploration assumption, with an explicit finite-horizon bound of $-\log w(p^\star)$ on the expected cumulative predictive KL.
\item Fixed-false-program elimination (Lemma): any fixed false program loses posterior odds almost surely under cumulative R\'enyi-$\tfrac12$ separation, via an exact nonnegative-martingale argument.
\item Semantic consistency (Theorem): the posterior concentrates almost surely on the semantic equivalence class $[p^\star]$ of the true generator under cumulative aggregate separation of the posterior false-mixture predictive law from the true kernel, proved by the same martingale transform applied to the full false-complement posterior odds.
\item Uniform exploratory access (Proposition): an explicit sufficient condition for the aggregate-separation hypothesis, stated in terms of a history-measurable separating action with a nonvanishing probability floor under the operating policy.
\item $\epsilon$-mixed expected-free-energy corollary: an expected-free-energy policy, mixed with an exploratory component of arbitrarily small weight $\epsilon>0$, inherits the uniform-access hypothesis and therefore the semantic-consistency conclusion, provided the same structural separation premise holds.
\end{enumerate}

The open asymptotic question is no longer whether interactive semantic consistency can be proved. It is whether the policy-side separation premise can be weakened and derived from more intrinsic properties of the exact expected-free-energy control rule. We flag this as open rather than finessing it.

The abstract, introduction, positioning statement, Section~8 opener, and discussion have been revised so that the paper neither understates nor overstates what is now proved.

\section*{Response to RMC2: No Empirical Content (Still)}

We agree with the reviewer's framing and do not try to argue the paper into a different venue. The present submission is a theory paper with an evaluation protocol, and it is deliberately silent on empirical results. The companion empirical paper is in preparation, and the three-level framework in this paper was organized explicitly so that the ideal target, the proxy measurements, and the mechanistic-challenge tolerances are stable objects against which that subsequent empirical work can be reported.

For the present round we do not ask the reviewer to treat an empirical gap as closed. We ask only that the theory be evaluated on its own terms: finite-history identities, asymptotic theorems in the interactive realizable case, and a specification (not a demonstration) of the mechanistic compliance protocol.

\section*{Response to RMC3: Synthesis-versus-Contribution}

We accept the reviewer's premise that the relevant question is what the assembly buys that the components do not already provide separately.

The synthesis-specific technical payoff is the Section~8 theorem package. Once hidden causes are interactive executable programs under a universal code prior, the inferential target is no longer a generic latent parameter but the semantic equivalence class $[p^\star]$ of the true explanatory program. Because actions are endogenous, distinguishability depends on which interventions the policy actually visits, and the relevant separation must be measured against a posterior false-mixture predictive law rather than against any fixed alternative. The revision proves semantic-class concentration in exactly that regime. This is not a direct corollary of:
\begin{enumerate}
\item Ortega--Braun, who provide the variational bounded-rationality form for generic reference priors;
\item Lindley / Chaloner--Verdinelli / MacKay, who provide information-gain for latent parameters in generic experimental-design settings;
\item Schmidhuber, who provides a compression-curiosity reading at the level of intrinsic reward; or
\item Hutter / Leike, who provide the universal-prior interactive-semimeasure apparatus but not this variational free-energy framing or a semantic-class identification theorem in this form.
\end{enumerate}

What the synthesis adds technically is therefore not another derivation of information gain but a semantic-class recovery theorem for interactive explanation programs. The ideal / behavioral proxy / mechanistic challenge three-level framework is the companion conceptual contribution: it explains how that ideal target could be evaluated in practice despite incomputability and limited interpretability.

The prose around the deterministic-discovery corollary, the related-work section, and the discussion has been revised so that those passages no longer suggest the contribution is merely a Solomonoff-flavored restatement of Lindley or Schmidhuber.

\section*{Response to RMC4: The Mechanistic Challenge Remains Unexercised}

We accept the reviewer's observation and have not tried to hide it. The mechanistic challenge in this paper is a specification rather than a run. The revision strengthens the specification in three ways: (a) by removing the underdeveloped hard-to-vary third test (see RMC5), so that the protocol is now a two-quantity joint test rather than a three-test hybrid; (b) by sharpening the interpretability-substrate subsection to say, for each cited tooling program, precisely what kind of substrate it supplies toward exposing $q_t$ or $\mathcal{G}_t$; and (c) by making compositional compliance a two-component statement that matches the revised protocol.

Running the protocol on a model is an empirical task, not a theoretical one. We agree that calibration of the protocol on at least one real system is a necessary next step; we have scheduled that exercise alongside the companion empirical paper and note it as such in the discussion.

\section*{Response to RMC5: Hard-to-Vary Criterion Under-Developed}

We agree with the reviewer's diagnosis and have removed the hard-to-vary material from the paper. Specifically:
\begin{enumerate}
\item The list of explanatory criteria in Section~2 is now three items rather than four, and the criteria-level table reflects the change.
\item The proxy map table in Section~7 no longer includes a ``hard-to-vary structure'' row.
\item The mechanistic-challenge compliance protocol in Section~8 no longer contains the perturbation test (iii); the specification is now a joint two-quantity test on $q_t$ and $\mathcal{G}_t$, and the statistical-power statement, the requirements table, and the compositional-compliance subsection have been updated to match.
\item The ``Hard-to-vary structure'' paragraph in Section~11 (Discussion) has been deleted, and the Deutsch (2011) citation has been removed from the bibliography.
\end{enumerate}

We agree with the reviewer that the criterion, while philosophically attractive, was not supported by a theorem connecting the universal prior to program-neighborhood behavior, and the DSL-edit perturbation instantiation was a modeling choice whose calibration had not been done. Removing the criterion narrows the paper's claims to the ones the theorems and the compliance protocol actually support. The removal also tightens the mechanistic-challenge specification: a joint test of $q_t$-compliance and $\mathcal{G}_t$-ranking compliance is a single-system assessment in a way that a three-test hybrid was not.

\section*{Response to N1: Rhetorical Slippage Between ``Variationally Specified'' and ``Normatively Complete''}

The revision uses ``variationally specified'' consistently throughout the abstract, introduction, positioning statement, Section~8, and discussion. The only remaining uses of ``principled'' occur where the paper is explicitly discussing whether the substitution of an algorithmic prior for a stipulated domain prior counts as principled at all---a question the paper now frames as a philosophical choice one can reasonably accept or question. The conclusion has been rewritten to align with the reference-machine caveat and no longer implies that a prior-justification problem has been settled.

\section*{Response to N2: Reference-Machine Invariance Partly Concedes Its Own Motivation}

We agree with the reviewer that the invariance discussion forces a narrower and more careful motivational claim than earlier framing suggested. The strongest defensible claim is:
\begin{enumerate}
\item the algorithmic prior replaces many domain-local prior choices with one global coding convention,
\item that convention has a precise universality-up-to-constant property across representable hypotheses, and
\item once fixed, it yields an exact variational identity, an exact connection to description length, a specializable expected-free-energy decomposition, and an interactive semantic-recovery theorem under explicit separation and exploration conditions.
\end{enumerate}

We do not claim finite-sample uniqueness or objectivity of the prior. The revised abstract now states explicitly that the universal measure is fixed relative to a chosen universal prefix machine and is therefore unique only up to an additive invariance constant. The conclusion has been rewritten in the same spirit: the machine is a global modeling choice rather than an eliminated stipulation, and the invariance is asymptotic in nature. Section~8.1 continues to say that the invariance constant can be arbitrarily large on any fixed finite-sample problem.

\section*{Response to N3: Recent Interpretability Citations Are Load-Bearing But Thin}

We agree. The interpretability-substrate subsection (Section~8.3) has been rewritten so that each cited line carries a concrete statement of what it actually supplies toward the compliance protocol. Linear probes (Alain--Bengio; Belinkov--Glass) are trained readouts from frozen internal representations; extending them to produce a calibrated distribution over a fixed finite program set is a parameterization question, not a new research problem. Mechanistic-circuit work (Olsson et al.; Wang et al.) identifies named subroutines whose activity is measurable across inputs and supplies a substrate that could be mapped to posterior weight on corresponding candidate programs. Sparse-autoencoder work (Cunningham et al.; Templeton et al.) decomposes hidden states into sparse interpretable feature directions and gives, where a candidate program set factors along those features, a principled starting point for estimating $q_t$. Targeted parameter edits such as ROME (Meng et al.) and activation-patching techniques supply the causal-manipulation tooling that would let an evaluator check whether an exposed $q_t$ or $\mathcal{G}_t$-ranking is dispositional rather than a post-hoc readout.

None of these tools, taken off the shelf, currently exposes $q_t$ or $\mathcal{G}_t$ in the compliance-protocol form, but each supplies a specific substrate that a compliant system would combine.

\section*{Response to N4: Ortega--Braun Comparison Is Weak in One Specific Respect}

We agree. The Delta paragraph in the related-work cluster on information-theoretic bounded rationality now states the identification explicitly. The Ortega--Braun functional has the form $\mathbb{E}_q[L] + \tau\,\mathrm{KL}(q\|q_0)$ for a reference policy or prior $q_0$ and a cost coefficient $\tau$. The scaled functional $\mathcal{F}_t^\tau$ of Section~5 is of exactly that form under the identification $q_0 \leftrightarrow w$ and the same coefficient $\tau$; the free-energy minimizer $q_t^{\star,\tau}$ is the corresponding Gibbs solution. The substantive difference is therefore not the variational form itself but the specific reference: Ortega--Braun typically take $q_0$ as a problem-specific reference policy or domain prior, while we take it as the universal prior over self-delimiting interactive programs, fixed once per choice of universal prefix machine. We state this identification in the Delta paragraph rather than treating Ortega--Braun as an orthogonal component.

\section*{Summary of Changes to the Manuscript}

\begin{itemize}
\item Abstract: reference-machine caveat added.
\item Section~2 (Motivation): hard-to-vary criterion removed; four-criterion list now three; criteria-level table updated.
\item Section~3 (Related Work): Ortega--Braun Delta paragraph sharpened to state the functional identification explicitly.
\item Sections~4--6: unchanged in content.
\item Section~7 (Behavioral Proxy): hard-to-vary row removed from the proxy-map table.
\item Section~8 (Reference-Machine Invariance, Posterior Consistency, and Exploration): predictive merging, fixed-false-program elimination, semantic consistency, exploration sufficiency, and the $\epsilon$-mixed expected-free-energy corollary are all proved. Retained in full from the previous revision.
\item Section~9 (Mechanistic Challenge): reduced from three tests to two, on $q_t$ and $\mathcal{G}_t$. Statistical-power, requirements-table, and compositional-compliance subsections updated. Interpretability-substrate subsection rewritten with one concrete sentence per cited line.
\item Section~11 (Discussion): hard-to-vary paragraph removed.
\item Conclusion: rewritten to align with the reference-machine caveat (``global modeling choice rather than an eliminated stipulation''; ``unique only up to an additive invariance constant that is asymptotic in nature''), and to reflect the now-two-quantity compliance protocol.
\item Bibliography: Deutsch (2011) entry removed. All other entries retained.
\end{itemize}

\section*{What Remains Open}

The revision does not close every open question; it closes several and narrows others to statements the paper can honestly sustain. The remaining open items, in order of importance, are:
\begin{enumerate}
\item Empirical validation of the behavioral proxy on a structured rule-recovery benchmark (deferred to the companion paper).
\item A demonstration of the mechanistic-challenge compliance protocol on at least one real system, for example a trained DSL interpreter or a DreamCoder-style agent, to calibrate the tolerances $\varepsilon_q$ and $\tau_\mathcal{G}$.
\item Weakening the policy-side separation premise in Section~8 and deriving it from more intrinsic properties of the exact expected-free-energy control rule.
\end{enumerate}

Items 1 and 2 are empirical and out of scope for this submission. Item 3 is a genuine theoretical next step; we flag it rather than obscure it.

\end{document}
